% Generated by scripts/gen_blueprint.py from blueprint/nodes/01-introduction.toml. Do not edit.

\chapter{Introduction}
\label{chap:01-introduction}

This blueprint covers the classical part of the Lean 4 library \texttt{diaz-modulus-lean}: the theorems of transcendental number theory in the table ``The classical theorems'' of its README, and every result their proofs use. They are the theorems of Hermite--Lindemann and Lindemann--Weierstrass, the Gelfond--Schneider theorem and the transcendence of $e^{\pi}$, the six exponentials theorem, the four exponentials theorem in transcendence degree one, Waldschmidt's theorem of 1973 with Schneider's eighth problem, and Roy's lemma on singular spaces of matrices; on the way, Waldschmidt's transcendence criterion and his zero count for exponential polynomials. No new transcendence theorem is claimed. The library's results on Diaz's conjecture are not part of this blueprint.

Every result here is a separate statement with its own Lean proof, and every one is proved: the library depends only on Lean's three standard axioms. Each is a node proved on the platform Prove2Me (\url{https://prove2.me}) and ported to the library, and the dependencies shown are the imports of the accepted proofs, with one exception noted at Theorem~\ref{thm:hermite_lindemann}. Each statement names its Lean declaration and was checked against it; where a Prove2Me page states a result in another form, the statement here follows the Lean declaration. An algebraic number is a complex number algebraic over $\Q$, a logarithm of an algebraic number is a complex number $\ell$ with $e^{\ell}$ algebraic, and $\N$ contains $0$.

Three parts come from elsewhere. The Lindemann--Weierstrass development is ported from Mathlib pull request \#28013 by Yuyang Zhao, which is not merged. The proof of the Gelfond--Schneider theorem restructures the formalisation of M.~Karatarakis and F.~Wiedijk \cite{KW26}. The statement \texttt{e\_pi\_transcendence} is a node posed on Prove2Me by another contributor; the proof here applies Gelfond--Schneider.

\textbf{[Origin and AI use: text to be agreed with Carlo.]}
